\chapter{Reproducibility Details}
\label{app:repro}
% 새 부록(방법 리뷰 P1-5, P1-6, P3-2). 근거는 각 절의 주석.

Settings, selection grids, eligibility outcomes, and resampling details
are given below as the stored contracts and reports state them; where a
detail is not documented there, the gap is marked. The model itself is
specified in Section~\ref{sec:method-baseline}. This appendix is not a
complete reproduction package, and the in-house results of the thesis
are restatements of the stored round reports rather than experiments
that a reader can repeat from this text alone.

\section{Code, environment, and what the archive lacks}
\label{sec:repro-code}

A minimal executable package prepared on 23 November 2026 contains the
inference code, the fitted preprocessor, and one checkpoint, the C-fold
REL\_GLOBAL model of seed 127. It requires Python 3.10 or later and
PyTorch 2.5.1 on CPU, and it states that the material needed to
reproduce the whole study is outside its scope. The training, run and
scoring programs of the v3 round remain on the original experiment
machine and are not in the archive available to us. That work was not
a Git repository, so no commit exists; the seals record SHA-256 hashes
of files instead (Table~\ref{tab:repro-hashes}). Training ran on a host with 251~GiB of memory and
128 CPUs, with an 8~GiB memory limit per process; v3 predictions were
generated with two CPU threads. The exact Python and PyTorch versions
used during the study were not recorded.

Each round has a seal file; the SHA-256 values we could read back are
those of the v3 contract and \texttt{fit\_v3.py}, of the generality and
ADS-24 contracts, and of the two plan files of the 135-run round, and
the other entries of the seal files were not read for this thesis. The
original machine holds the training
seed plans, split and source manifests, run ledgers and eligibility
seals (for example \texttt{TRAIN\_SEED\_PLAN.json},
\texttt{SPLIT.json}, and \texttt{ELIGIBILITY\_SEAL.json}), as well as
the seed-level score files that the reports cite
(\texttt{SCORED\_RESULTS.json}, \texttt{SEED\_RESULTS.csv}). Their
contents are not in the archive available to us. The seeds, generator
seeds, exclusion reasons and seed-level counts given in this appendix
and in Chapter~\ref{ch:results} were read or computed from these files
on the original machine, without new training or inference.

A reanalysis on the original machine on 25 November 2026 read the saved
probability files of every complete 30-seed group, 10,260 files in all,
and recomputed 355 group scores of the stored reports, among them every
group mean of the v3 round, with a largest difference of
$1.8\times10^{-15}$. Its per-run, per-seed table reproduces the v3 seed
standard deviations, worst seeds, and per-seed counts of
Chapter~\ref{ch:results}, and Figure~\ref{fig:v3-seeds} is drawn from
it. It did not recompute the criteria intervals at their sealed level;
the 95\% configuration intervals it computed for each collection have
the same signs. The probability files stay on the original machine.
% 출처: inputs/data/FINAL_CHECKS_20261125/UNIFIED_REANALYSIS_20261125/UNIFIED_REANALYSIS.md:3 (342 묶음, 10,260 파일, 355 점수, 최대 차이
%   1.78e−15); SAVED_SCORE_RECOMPUTE.csv (v3 PEKD4 15칸 포함); COLLECTION_COMPONENTS.csv (PEKD4 95% 구간).

What a reader can check is still limited. For the v3 round
(Table~\ref{tab:rounds}), the thesis does not include the training, run and scoring
programs, the scored row-level outputs, or the split and source
manifests. The minimal inference package, the training and scoring
code and the row-level scores are not published with this thesis. The
SHA-256 hashes, seed lists and bootstrap generator seeds given here
allow a stored result to be matched against the files it came from. A
reader can check that
Tables~\ref{tab:results-v3-nll} and~\ref{tab:results-v3-criteria}
transcribe the stored report, but cannot recompute the means, seed standard
deviations, worst seeds or intervals from row-level data without the
files on the original machine. The
sampling rules of the training batches are given in
Sections~\ref{sec:method-baseline} and~\ref{sec:method-variants}. The
thesis accordingly presents the v3 verdicts as recorded in that report
and confirmed by the later recomputation of its scores, not as an
experiment a reader can repeat from this text, and it draws no
conclusion from v3 that goes beyond what the report records.
% 7회차 방법 P1-1: 원자료를 보충자료로 낼 수 없으므로(atlas에 없음), v3를 "저장 보고서에 기록된
%   역사적 결과"로 한정했다. 06_results, 01_introduction, 99_conclusion에 같은 한정을 넣었다.
% 저자 결정(2026-11-24, 저자 확인 목록 5번): 최소 실행본·학습·채점 코드·행 단위 점수는 공개하지 않는다.
%   재생성 명령은 넣지 않는다. '요청 시 제공' 같은 약속 문구도 쓰지 않는다.
% 출처: sources/atlas_answers/repro.md 5번.
%   atlas OT/PROPHET_COPYRIGHT_SUBMISSION_20261123/package/README.md:3,7,20,23;
%   OT/PROPHET_ENV_CONSENSUS_KD_V3_20261120/PAPER_DEFINITIONS.md:16,28 (fit_v3.py 등 이름만);
%   OT/REL_GROUP_OBJECTIVES_FIXED_20261114/RUN_STATE.md:9 (251 GiB, 128 CPU);
%   OT/PROPHET_ENV_CONSENSUS_KD_V3_20261120/REPRO_FIX_LOG.md:3 (torch.set_num_threads(2));
%   OT/PROPHET_REL_GLOBAL_SEED_POPULATION_20261118/FINAL_REPORT_KO.md:13 (해시 봉인);
%   PROPHET_RESEARCH_TRIAL_ATLAS/FILE_INVENTORY.csv.gz (원본에만 있는 파일).
%   8 GiB 상한: repro.md 5번(학습 호스트).
% 출처(commit 없음): 원래 PC 의 작업 폴더와 /home/yonggamman/prophet-e4 에서 git rev-parse HEAD 가
%   "not a git repository"; 봉인 파일에도 commit 항목 없음(codex_original_machine.md 저장값 2–4).
% (해결, 저자 확인 5번) 최소 실행본(seed 127 checkpoint)과 코드는 공개하지 않는다(저자 결정 2026-11-24).

\begin{table}[htbp]
\centering
\caption[SHA-256 hashes recorded in the seals]{SHA-256 hashes recorded in the seals. The 135-run development
set seals its two plan files rather than one contract file.}
\label{tab:repro-hashes}
\scriptsize
\begin{tabular}{@{}lll@{}}
\toprule
Round & File & SHA-256 \\
\midrule
v3 & \texttt{CONTRACT\_V3.md} & \texttt{\begin{tabular}[t]{@{}l@{}}26cdaa793fc8074f90329229d2c5304d\\5f51145c135c9e5f6b56ccfc1200f56d\end{tabular}} \\
v3 & \texttt{fit\_v3.py} & \texttt{\begin{tabular}[t]{@{}l@{}}e5c51cd6ccb2b6f778f5e8c9b0f0e04a\\a5fb1d5326f2fedad45615afc526adc3\end{tabular}} \\
PAG1 & \texttt{ARCH\_GENERALITY\_CONTRACT\_SEALED.md} & \texttt{\begin{tabular}[t]{@{}l@{}}a2a7f3ff8071c0e1f15a3e240d6e843d\\56b84ead17b33e4846c0eb14fdf94b4d\end{tabular}} \\
ADS-24 & \texttt{ADS24\_CONTRACT\_SEALED.md} & \texttt{\begin{tabular}[t]{@{}l@{}}f93d4ea6b9221ce38c39dc1f290dea18\\9206ac57764f4a44ebeb3bb0048b5af8\end{tabular}} \\
135-run set & \texttt{BG\_MAIN\_PLAN.json} & \texttt{\begin{tabular}[t]{@{}l@{}}5be722717e56710c25e178e9c83c8409\\3a92f1fa7ad3255e7aac8c18505f0929\end{tabular}} \\
135-run set & \texttt{ORIGINAL\_MAIN\_PLAN.json} & \texttt{\begin{tabular}[t]{@{}l@{}}fbddf9ede43483267a0faa468e8903f9\\8c02d321cb5e0bc2a607fea7e1458ca2\end{tabular}} \\
\bottomrule
\end{tabular}
\end{table}
% 출처: 원래 PC: PROPHET_ENV_CONSENSUS_KD_V3_20261120/SEAL.json:1153 (v3 두 행);
%   원래 PC: ARCH_GENERALITY_STAGE1_20261122/SEAL.json:701 (PAG1); 원래 PC: ADS24_STAGE1/SEAL.json:10 (ADS-24);
%   원래 PC: 135실행 차수의 SEAL.json:1 (두 계획 파일; 계약 파일 하나의 hash는 없다).

\section{Training settings and seeds}
\label{sec:repro-training}

Every in-house fit of REL\_GLOBAL and of the CD students built on it
(Table~\ref{tab:variants}) uses the
603 runs of the C-fold FIT set with execution-weighted 18-way
cross-entropy, the representation-gradient penalty $\mathcal{R}_{\mathrm{GLOBAL}}$
with coefficient 0.1, and LogitNorm with $\tau=0.1$. Each fit runs
1,346 updates of 32 FIT occurrences with AdamW (weight decay 0.01,
gradient clipping at 2) and a cosine learning rate from 0.002 to
0.0002, and the last checkpoint is kept. Within a seed, all groups of a
round share the initial tensors and the FIT occurrence and delivery
trace. All in-house fits use the 404-only preprocessor of
Section~\ref{sec:method-baseline}. The three general encoders and
their CD students drop LogitNorm, $\mathcal{R}_{\mathrm{GLOBAL}}$, and
the relational residual; their settings are given in
Section~\ref{sec:repro-encoders}.
% 8회차 방법 P2-1: 범위를 REL_GLOBAL 계열로 한정하고 일반 인코더는 E.3절로 가리켰다.

The 30 REL\_GLOBAL seeds of the teacher population are 3674, 12050, 15301, 16635, 16893, 20436, 28035,
28756, 29771, 33674, 36720, 37406, 43142, 48192, 49208, 50572, 50836,
61616, 64832, 68498, 69205, 71526, 71556, 76939, 79043, 90835, 91513,
95246, 97386, and 99609. The 30 seeds of the v2 round were drawn without
replacement from 1,000,000--1,999,999 with PCG64 seeded by 2026111909
and sorted; the values themselves are recorded only in the round's seed
plan, which is not in the archive. The 30 new seeds of the v3 round are
2004673, 2067898, 2069477, 2077173, 2112904, 2122696, 2331561, 2339652,
2356542, 2361613, 2363065, 2400243, 2422305, 2446238, 2474153, 2482862,
2515409, 2574712, 2587606, 2613239, 2627292, 2672261, 2680405, 2751475,
2764668, 2855487, 2874630, 2877233, 2898425, and 2958656. The ADS-24 baseline seeds are 301001 to 301030. The five in-house teachers are the REL\_GLOBAL seeds 50572,
91513, 71556, 68498, and 37406; the five ADS-24 teachers are the
REL\_GLOBAL-ADS seeds 301030, 301026, 301025, 301003, and 301022. The
generality study and the method comparison use the seeds 4500001 to
4500030 for every group. In the FIT-input distillation group of the
method comparison, the one row per occurrence that receives the KL term
is drawn in advance with PCG64 seeded by 2026112305 plus the training
seed. Four-checkpoint weight averaging takes the mean of the weights at
updates 1,010, 1,122, 1,234, and 1,346 of the cross-entropy trajectory,
under the unchanged cosine schedule. The added comparisons of the final
checks reuse the seeds 4500001 to 4500030, the initial tensors, and the
FIT traces of the plain GRU and Transformer encoder; their random
teacher sets and input subsets were drawn before training and are
recorded in the stage's contract, which is not in the material
available to us.
% 근거: PAPER_DEFINITIONS.md §2 (C-fold 603 FIT, 404-only 전처리, 1,346×32, 같은 초기
% tensor·FIT trace, AdamW .002→.0002 cosine, wd .01, clip 2, 마지막 checkpoint; 상위 5
% seed); ADS24_STAGE2_REPORT.md "계보" (ADS 교사); ARCH_GENERALITY_CONTRACT_SEALED.md
% (4500001–4500030); METHOD_BASELINES_CONTRACT_SEALED.md (FIT_KD PCG64 2026112305+seed,
% SWA update 1010·1122·1234·1346).
% 출처(seed 목록): OT/PROPHET_INDEPENDENT_CONFIRM_CONTRACT_DRAFT_20261118/CONTRACT_DRAFT_KO.md:15 (30 seed);
%   OT/PROPHET_ENV_CONSENSUS_KD_V2_20261119/CONTRACT_V2.md:21 (PCG64 2026111909);
%   OT/PROPHET_ENV_CONSENSUS_KD_V3_20261120/ADS24_CONTRACT_SEALED.md:33 (301001–301030);
%   OT/PROPHET_ENV_CONSENSUS_KD_V3_20261120/RUN_STATE.md:4 (v3 새 30 seed).
% 출처(v3 30 seed 값): 원래 PC: PROPHET_ENV_CONSENSUS_KD_V3_20261120/TRAIN_SEED_PLAN.json:3.

\section{General encoders}
\label{sec:repro-encoders}

Like REL\_GLOBAL, the three general encoders read the receipt-bounded
prefix, with the 48-dimensional event representation and the age
and receipt features, and map an empty prefix to a learned common
state. The LSTM encoder is one unidirectional layer whose last state
feeds a softmax head. The Transformer encoder has two layers and four
attention heads and reads a classification token. The plain GRU is one
unidirectional layer whose last state feeds the head. All three train
with execution-weighted cross-entropy, AdamW with weight decay 0.01, a
cosine schedule down to 10\% of the initial learning rate, gradient
clipping at 2, 1,346 updates of 32 occurrences, and the last
checkpoint. On ADS-24 the alert window is 512.

Width and learning rate were chosen once per encoder and dataset from
four candidates, width $\{64, 128\}$ crossed with two learning rates,
$\{0.001, 0.002\}$ for the LSTM and GRU and $\{0.0005, 0.001\}$ for the
Transformer, each trained for 1,346 updates with a pilot seed outside
the contract. The selection score was the execution-weighted NLL on a
holdout inside FIT: 40 training and 8 held-out runs stratified by
scenario on ADS-24, and 483 and 120 runs split by a hash of the run
identifier family in-house. Ties would have gone to the smaller width
and then the lower learning rate. Table~\ref{tab:repro-encoders} lists
the selected settings; the method comparison reuses those of the plain
GRU. Within an encoder and dataset, the base model and the CD student
share the selected setting, so each contrast of
Chapter~\ref{ch:generality} compares two models of the same setting.
Comparisons between encoders are less secure: each setting was chosen
with one pilot seed, the selection was not repeated across seeds, and
the absolute ranking of the encoders is conditional on that choice.

\begin{table}[htbp]
\centering
\caption[Settings selected for the general encoders]{Settings selected for the general encoders by FIT-internal
validation (pilot seed outside the contract).}
\label{tab:repro-encoders}
\small
\begin{tabular}{@{}llrrrr@{}}
\toprule
Data & Encoder & Width & Learning rate & Holdout NLL & Parameters \\
\midrule
ADS-24   & LSTM (DeepLog-style)        & 64  & 0.002 & 1.541 & 46,213 \\
ADS-24   & Transformer (LogBERT-style) & 128 & 0.001 & 1.497 & 417,029 \\
ADS-24   & plain GRU                   & 128 & 0.001 & 1.498 & 119,301 \\
In-house & LSTM (DeepLog-style)        & 128 & 0.002 & 0.099 & 155,834 \\
In-house & Transformer (LogBERT-style) & 64  & 0.001 & 0.140 & 116,538 \\
In-house & plain GRU                   & 128 & 0.002 & 0.112 & 122,810 \\
\bottomrule
\end{tabular}
\end{table}
% 근거: ARCH_GENERALITY_CONTRACT_SEALED.md "세 적용판", "FIT 전용 선택 및 1단계"
% (후보 격자, ADS 40/8 시나리오 층화, PROPHET 483/120 실행 ID 계열 해시 보류,
% 내부 보류 실행 동등 NLL 최소, 동률 규칙, pilot seed 3300001, 선택 표의 내부 NLL
% 1.540549/1.497478/1.497624/0.099077/0.139591/0.112230과 활성 파라미터 수);
% METHOD_BASELINES_CONTRACT_SEALED.md (방법 비교의 PLAIN_GRU 폭·학습률).

\section{Eligibility, exclusions, and safety stops}
\label{sec:repro-eligibility}

Table~\ref{tab:repro-eligibility} gives, for each in-house collection
that the main chapters score, the planned and attempted executions, the
eligible set, and the recorded reasons for the rest. Preparation runs
were separate identifiers in every round and never entered an
evaluation. The safety-stop rules of the contracts read only the state
of the collection. The v2 contract stopped on 256~MiB of free disk,
8~GiB of child memory, 420~s per attempt, an uncertain cleanup, or three
consecutive source failures. The PAG1 contract kept these and added an
18-hour limit and a source completion rate below 0.85 after 20
attempts, and the method comparison applied the rule of three
consecutive source failures.

\begin{table}[htbp]
\centering
\caption[Eligibility and stopping per in-house collection]{Eligibility and stopping per in-house collection. Each planned
identifier was attempted at most once, and nothing was replaced.}
\label{tab:repro-eligibility}
\small
\begin{tabularx}{\textwidth}{@{}>{\raggedright\arraybackslash}p{3.1cm}rrr>{\raggedright\arraybackslash}X@{}}
\toprule
Collection & Planned & Attempted & \shortstack[r]{Eligible\\runs / rows} & Not eligible or not attempted \\
\midrule
\texttt{OBSERVED\_HEAD} check (later the 10-seed round) & 160 & 160 & 157 / 1,611
  & three background runs whose COPY START, DISCOVERY START, or DISCOVERY
    STOP record could not be bound to its source request \\
v2 round, first & 160 & 117 & 116 / 1,187
  & stopped as an uncertain cleanup after a network fixture listener was
    reported as exited; 43 identifiers not attempted \\
v2 round, recheck & 160 & 160 & 159 / 1,633
  & one background run whose source START at the last boundary could not
    be bound \\
v3 round & 160 & 160 & 157 / 1,611
  & three background runs that failed during source collection: a
    required SOURCE\_ACTION\_STOP missing (\texttt{PEKD4B0006}), a required
    SOURCE\_ACTION\_START missing (\texttt{PEKD4B0087}), and a source
    service-health timeout (\texttt{PEKD4B0112}) \\
PAG1 & 160 & 160 & 156 / 1,600
  & four background runs whose source action START (\texttt{PAG1B0023},
    \texttt{PAG1B0037}, \texttt{PAG1B0061}) or STOP (\texttt{PAG1B0097})
    record was not found$^{*}$ \\
Method comparison & 160 & 119 & 110 / 1,119
  & five source failures; four runs with an unbound source START or STOP;
    stopped after three consecutive source failures (agent or forwarder
    not ready); 41 identifiers not attempted \\
Final checks, first (PFC1) & --- & --- & 0
  & safety stop; kept in the ledger, counts not in the material available to us \\
Final checks (PFC2) & 160 & 160 & 155 / 1,589
  & five background runs whose source START (\texttt{PFC2B0023},
    \texttt{PFC2B0046}, \texttt{PFC2B0060}) or STOP (\texttt{PFC2B0005},
    \texttt{PFC2B0113}) record was not found$^{*}$ \\
\bottomrule
\end{tabularx}
\par\smallskip
\begin{minipage}{\textwidth}\footnotesize\setstretch{1}
$^{*}$Reasons as recorded in the sealed eligibility file. The four PAG1,
the five PFC2, and the three v3 exclusions are all failures at the
source collection stage, before any prediction was made, and none
depends on a score.
\end{minipage}
\end{table}
% 근거: OBSERVED_HEAD_STAGE1_CONFIRM FINAL_REPORT_KO.md:11 (POHCB0005 COPY START,
% POHCB0014 DISCOVERY START, POHCB0079 DISCOVERY STOP 미결속); PAPER_DEFINITIONS.md §3 및
% PROPHET_ENV_CONSENSUS_KD_V2_20261119 (PEKD2O0011 NETWORK_FIXTURE_LISTENER_EXITED 후
% TIMEOUT_OR_CLEANUP_UNCERTAIN 중지, 43 미시도); V2_DIAG_RECHECK FINAL_REPORT_KO.md:5
% (PEKD3B0089); 출처(v3 제외 사유): atlas OT/PROPHET_ENV_CONSENSUS_KD_V3_20261120/RUN_STATE.md:7
%   (세 ID와 사유, 보충 없음; repro.md 8번);
% ARCH_GENERALITY_STAGE1_REPORT.md (배경 124/128 적격); METHOD_BASELINES_REPORT.md
% "실행과 재현" (SOURCE_FAILED 5, SOURCE_ACTION_START_NOT_FOUND 3, STOP_NOT_FOUND 1,
% NOT_ATTEMPTED 41) 및 TECHNICAL_LEDGER.md (AGENT_OR_FORWARDER_NOT_READY 세 번 연속).
% 안전 중지 규칙: CONTRACT_V2.md:15; ARCH_GENERALITY_CONTRACT_SEALED.md "새 PROPHET 수집".
% 출처(PAG1 제외): 원래 PC: ARCH_GENERALITY_STAGE1_20261122/ELIGIBILITY_SEAL.json:19
%   (ValueError: SOURCE_ACTION_START_NOT_FOUND:11/10/9, SOURCE_ACTION_STOP_NOT_FOUND:12).

\section{Resampling}
\label{sec:repro-resampling}

Every contrast interval comes from 100,000 resamples drawn with PCG64;
the stage-1 gate of the ADS-24 study used 10,000 resamples of
executions within scenarios. In-house,
a resample draws complete configurations and moves their repetitions
together, and the two A+C layouts are resampled separately and then
averaged 1:1 where the score combines them; on ADS-24, a resample draws
executions within scenarios. The fixed seeds of the generators are
2026112203, 2026112204, and 2026112205 for the first stage of the
generality study (ADS-24, in-house, and the combined A+C score),
2026112206, 2026112207, and 2026112208 for its second stage and the
PAG1 re-inference, 2026112303 and 2026112304 for the method
comparison on ADS-24 and in-house, 2026112314 and 2026112315 for the
added comparisons of the final checks on ADS-24 and in-house,
2026112501 for the post hoc pooling of collections, and 2026112003 for
the second stage of Chapter~\ref{ch:ads24}, where the three main
contrasts share one set of indices. The v3 contract specifies 100,000 configuration-paired
resamples and a Bonferroni level of 98.3333\%; its scoring code seeds
the generator with 2026112021 for \texttt{A+C FLAT}, 2026112022 for
\texttt{A+C NESTED}, and 2026112023 for \texttt{ORIGINAL\_18}.
% 근거: ARCH_GENERALITY_CONTRACT_SEALED.md "blind 채점과 2단계 관문";
% ARCH_GENERALITY_STAGE2_PREFIT_SEAL.md; METHOD_BASELINES_CONTRACT_SEALED.md "주 평가와 결정".
% 출처: atlas OT/PROPHET_ENV_CONSENSUS_KD_V3_20261120/ADS24_STAGE2_REPORT.md:15 (PCG64 2026112003);
%   CONTRACT_V3.md:11 (v3) (repro.md 9번).
% 출처(v3 생성기 seed): 원래 PC: PROPHET_ENV_CONSENSUS_KD_V3_20261120/score_v3.py:33, SEAL.json:1104
%   (numpy PCG64, 100,000회; 보조 C_ONLY 2026112024, B_PLUS_C 2026112025는 본문에 넣지 않았다).

\section{Point estimates on complete configurations}
\label{sec:repro-complete}

In the generality study, in-house point estimates use every eligible
run, whereas intervals use the configurations whose repetitions are all
eligible, so the two summarize slightly different populations.
Table~\ref{tab:repro-complete} gives both point estimates for the A+C
layouts of the generality study. The two differ by at most 0.011, and
the direction of every contrast is the same under either. In the v3
round the three exclusions are background runs, so at most three of the
64 background configurations (16 per condition, two repetitions each)
lost a repetition, and the seen-condition configurations are all
complete. Its three criteria rest on 15, 16 and 8 complete
configurations; their complete-only estimates equal the stored ones.
Leaving out any one configuration kept the sign of all three: the
estimates ranged over $[-0.049, -0.041]$, $[-0.054, -0.047]$ and
$[-0.075, -0.039]$.
% 출처: 원래 PC: PROPHET_ENV_CONSENSUS_KD_V3_20261120/SCORED_RESULTS.json:1, :1191; score_v3.py:84;
%   evaluate_v3.py:134 (codex_original_machine.md 저장값 5). 구성별 LOO 값은 원 답 파일에만 둔다.

\begin{table}[htbp]
\centering
\caption[Mean NLL on the PAG1 runs, all eligible runs against complete configurations]{Mean NLL over 30 seeds on the PAG1 runs: all eligible runs
against complete configurations only (\texttt{A+C FLAT}: 15
configurations, 30 runs; \texttt{A+C NESTED}: 14 configurations, 28
runs).}
\label{tab:repro-complete}
\small
\begin{tabular}{@{}llrrrr@{}}
\toprule
& & \multicolumn{2}{c}{Base model} & \multicolumn{2}{c}{CD student} \\
\cmidrule(lr){3-4}\cmidrule(l){5-6}
Encoder & Layout & All eligible & Complete & All eligible & Complete \\
\midrule
LSTM (DeepLog-style)        & FLAT   & 0.651 & 0.654 & 0.486 & 0.490 \\
LSTM (DeepLog-style)        & NESTED & 0.638 & 0.640 & 0.465 & 0.463 \\
Transformer (LogBERT-style) & FLAT   & 1.050 & 1.058 & 0.548 & 0.558 \\
Transformer (LogBERT-style) & NESTED & 1.041 & 1.034 & 0.520 & 0.514 \\
plain GRU                   & FLAT   & 1.139 & 1.143 & 0.258 & 0.261 \\
plain GRU                   & NESTED & 1.122 & 1.127 & 0.249 & 0.247 \\
\bottomrule
\end{tabular}
\end{table}
% 근거: ARCH_GENERALITY_STAGE2_REPORT.md "PROPHET 전체 적격 실행과 완전 구성 점추정"
% (최대 차이 Transformer FLAT KD 0.547842 대 0.558368, 0.0105).
